\documentclass[10pt,conference]{IEEEtran}

% ===================== PACKAGES =====================
\usepackage[utf8]{inputenc}
\usepackage[T1]{fontenc}
\usepackage{amsmath,amssymb,amsfonts,mathtools}
\usepackage{physics}
\usepackage[draft]{graphicx}   % NOTE: draft mode is the SAFE DEFAULT, so the
                                % document still compiles for any figure whose
                                % file is not yet present (e.g. future
                                % double-well figures). Figures that DO have a
                                % real file attached use the per-image
                                % override `draft=false` in their
                                % \includegraphics call, so they render
                                % normally regardless of the package default.
\graphicspath{{../../figures/}{../../figures/HO/}}
\usepackage{booktabs}
\usepackage{array}
\usepackage{url}
\usepackage[hidelinks]{hyperref}

\begin{document}

\title{A Modular Numerical Pipeline for the Quantum-to-Classical\\
Transition of Heat Capacity: DVR-Based Simulation\\
of Smooth 1-D Potentials}

\author{
\IEEEauthorblockN{Nadav Jean}
\IEEEauthorblockA{
\textit{Department of Chemical Engineering}\\
\textit{Technion -- Israel Institute of Technology}\\
Haifa, Israel\\
Course 00540367 -- Research Project 1, Spring 2026 (Semester 6)\\
Supervisor: Dr.\ David Gelbwaser-Klimovsky \quad Student ID: [insert]}
}

\maketitle

\begin{abstract}
A heat engine's cycle moves an amount of heat set directly by its working substance's heat capacity, so $C_v$ is, to first approximation, a proxy for how much heat -- and therefore how much power -- a cycle built from that substance can move. A classical working substance has a comparatively featureless $C_v(T)$. A quantum system with a small, effectively finite number of accessible energy levels \emph{can instead} show a Schottky anomaly: a sudden rise and fall of $C_v(T)$ as the restricted ladder of levels first activates and then saturates \cite{gelbwaser2018,kozin2025} -- a finite level count is necessary for this behavior, but far from every such system displays it. The long-term goal of this research line is to find a system where this anomaly lets $C_v$, and hence engine performance, transiently exceed its classical counterpart, in the spirit of recent quantum-thermodynamics results showing that enhanced energy fluctuations can let a heat engine approach the Carnot bound more closely at finite power \cite{campisifazio2016,ye2017}. This project builds the numerical groundwork for that search: a pipeline that computes $C_v(T)$, quantum and classical, for any smooth one-dimensional potential, using a Colbert--Miller sinc discrete variable representation (DVR) \cite{colbertmiller1992,lightcarrington2000} and an $\hbar$-scaling parameter $\xi$ that locates the classical limit \cite{gelbwaser2018}. Because most potentials of real interest have no closed-form solution, the pipeline validates itself without ever leaning on one: every spectrum is checked against an independently generated, higher-resolution numerical reference rather than an analytic formula, and a companion check rejects the artificial collapse of $C_v$ that can mimic a Schottky anomaly once a finite computed spectrum is exhausted. This closed-form-free methodology is itself certified, once, against the one system where the exact answer is independently known -- the harmonic oscillator (HO) -- where it reproduces the analytic result to machine precision, and is now being extended to a quartic double well.
\end{abstract}

\begin{IEEEkeywords}
quantum thermodynamics, discrete variable representation, heat capacity, classical limit, numerical validation, quantum heat engines
\end{IEEEkeywords}

% ===================== I. INTRODUCTION =====================
\section{Introduction}
\label{sec:intro}

\subsection{Motivation: The Schottky Anomaly as a Candidate Mechanism}
The broader aim of this project is to find operating regimes in which a quantum heat machine outperforms its classical counterpart -- not by exceeding the Carnot efficiency, which the second law forbids, but by \emph{approaching} that bound more closely while still delivering finite power. Because a heat engine's cycle moves an amount of heat set by its working substance's heat capacity, $C_v$ is, to first approximation, a proxy for how much heat -- and hence power -- a cycle can move. A classical substance's $C_v(T)$ is comparatively static; a quantum system with a small, effectively finite number of accessible levels can instead show a Schottky anomaly: at low $T$ only the ground state is populated ($C_v\to0$); $C_v$ rises sharply once $k_BT$ becomes comparable to the level spacing; and, because the available levels then saturate while further heating adds little entropy, $C_v$ falls again rather than settling onto a classical plateau \cite{gelbwaser2018,kozin2025}. A restricted, effectively finite level count is a \emph{necessary} ingredient for this behavior, but it is not sufficient -- most quantum systems with many levels do not display a Schottky peak at all -- so the anomaly is a candidate mechanism to be searched for, not something to be assumed.

Recent work has shown that a working substance with strongly enhanced energy fluctuations can let a heat engine approach the Carnot point without sacrificing power \cite{campisifazio2016}, and that the heat capacity of a small quantum system directly constrains the work a cyclic engine can extract \cite{ye2017}. Since $C_v(T)=k_B\beta^2\mathrm{Var}(E)$ is by definition a measure of those fluctuations, a genuine Schottky enhancement -- a system whose \emph{quantum} $C_v$ transiently exceeds its \emph{classical} value -- is exactly the kind of effect this project is searching for. Concretely, the goal is to find such a system, quantify how the enhancement would translate into engine efficiency relative to the classical limit, test whether it generalizes across a family of potentials, and use that to propose a concrete experimental realization. This report documents the groundwork -- a general-purpose numerical tool that computes $C_v(T)$, quantum and classical, for any smooth one-dimensional potential, without which a genuine enhancement could not be searched for systematically.

\subsection{Approach: A Self-Validating, Closed-Form-Free Pipeline}
Two ingredients make the pipeline usable on an arbitrary potential. First, the classical limit is reached not by taking $\hbar\to0$ directly -- $\hbar$ is a fixed constant of nature, not a quantity that can actually be dialed down -- but by scanning a dimensionless factor $\xi$ that plays the role of $1/\hbar$, realized instead by rescaling temperature and potential energy, quantities the calculation is genuinely free to vary; this is what makes the $\xi$-scan a physically implementable route to $\hbar\to0$ (Section~\ref{sec:theory}). Second -- and central to this report -- every energy spectrum the pipeline produces is checked not against an analytic formula, which does not exist for most potentials of interest, but against an independently generated, higher-resolution numerical reference computed by the same DVR solver (Section~\ref{sec:validation}). This self-consistent ground truth is the only one available once the project moves past the textbook cases, so this semester's central deliverable is showing that it can be trusted: Section~\ref{sec:validation} confirms it agrees with the HO, the one case where the exact answer is known, to machine precision.

% ===================== II. THEORY =====================
\section{Theoretical Background}
\label{sec:theory}
Throughout, the working substance is assumed to be in thermal equilibrium with a bath at temperature $T$: every $C_v(T)$ in this report is the equilibrium, canonical-ensemble heat capacity of a fixed spectrum $\{E_n\}$, not a finite-time or non-equilibrium quantity. This is the standard starting point for analyzing a quasi-static engine cycle built from a sequence of such equilibrium states. A real engine operating at the finite power discussed in Section~\ref{sec:intro} is necessarily not quasi-static, so it inherently involves genuinely non-equilibrium, finite-time corrections to this idealization; modeling those corrections directly is beyond this project's scope.

For a system with discrete energy levels $\{E_n\}$ at inverse temperature $\beta=1/(k_BT)$, the heat capacity follows from the variance of the energy,
\begin{equation}
C_v(T) = k_B \beta^2 \left[\expval{E^2} - \expval{E}^2\right] = k_B\beta^2\,\mathrm{Var}(E).
\label{eq:cv_general}
\end{equation}
This holds whether $\{E_n\}$ comes from a closed-form expression or a numerical diagonalization.

Following \cite{gelbwaser2018} (Eq.~S7 therein), the energy levels of a potential scaled by $\xi^2$ obey
\begin{equation}
\frac{E_n\!\left(\hbar_{\text{eff}}=\hbar,\ \xi^2 V(x)\right)}{\xi^2} = E_n\!\left(\hbar_{\text{eff}}=\frac{\hbar}{\xi},\ V(x)\right),
\label{eq:scaling}
\end{equation}
which is realized numerically by scaling both temperature and the spectrum by $\xi^2$ in Eq.~\eqref{eq:cv_general}, and is exactly equivalent to scaling $\hbar$ (which is a physical constant and so cannot realistically be scaled). Any computed spectrum is truncated at a finite highest level $E_{\max}$, so the plateau in $C_v(\xi)$ this produces can only be trusted inside a finite window: $\xi$ must be large enough for the level spacing to look continuous on the scale of $k_BT$, yet small enough that $E_{\max}$ stays thermally inaccessible,
\begin{equation}
\sqrt{\beta\,\Delta E} \;\ll\; \xi \;\ll\; \sqrt{\beta\,E_{\max}}.
\label{eq:xiwindow}
\end{equation}
Section~\ref{sec:validation} explains how the pipeline locates this window and recognizes when $\xi$ has left it.

For the HO, $E_n=\hbar\omega(n+\tfrac12)$ gives the closed-form Einstein heat capacity
\begin{equation}
C_v^{\text{HO}} = \frac{1}{4}k_B\left[\beta\hbar\omega\,\operatorname{csch}\!\left(\frac{\beta\hbar\omega}{2}\right)\right]^2 ,
\label{eq:cvho}
\end{equation}
tending to the equipartition value $k_B$ as $\beta\hbar\omega\to0$. Equation~\eqref{eq:cvho} is the closed-form answer the pipeline is checked against in Section~\ref{sec:validation}.

% ===================== III. DVR METHOD =====================
\section{The Discrete Variable Representation}
\label{sec:dvr}

\subsection{Why a Grid-Based Representation Works at All}
For potentials with no analytical spectrum, $\{E_n\}$ is obtained from a Colbert--Miller sinc discrete variable representation (DVR) \cite{colbertmiller1992,lightcarrington2000}, which represents the wavefunction directly by its values at a discrete set of evenly spaced grid points, rather than -- as in a traditional basis-set expansion -- by a set of coefficients multiplying smooth, closed-form functions (Hermite polynomials, plane waves, and the like). Because the grid is just a set of samples, it inherits a basic result from signal processing, the Nyquist--Shannon sampling theorem \cite{nyquist1928,shannon1949}: a signal can only be reconstructed faithfully from evenly spaced samples if it is sampled at least twice per cycle of its fastest-oscillating component; sampled any coarser, that fast-oscillating content does not simply disappear -- it reappears disguised as a spurious, slower-looking signal (\emph{aliasing}). For a wavefunction, the relevant ``frequency'' is the local wavenumber $k(x)$, which depends on both the state's energy $E$ and position $x$ (Section~\ref{sec:dvr}-D gives the specific value the grid spacing is set from): a state's local de Broglie wavelength $\lambda(x)=2\pi/k(x)$ sets how fast $\psi(x)$ oscillates at that point, so the grid spacing must resolve the \emph{shortest} such wavelength anywhere the wavefunction has significant amplitude, or that high-momentum content aliases back into the low-lying, physically relevant part of the computed spectrum as 'fake' eigenvalues.

\subsection{The Kinetic-Energy Matrix}
\label{sec:kinetic}
On a grid of $N$ points with spacing $\Delta x$, the potential is trivially diagonal, $V_{ij}=V(x_i)\delta_{ij}$; the entire method rests on the Colbert--Miller closed-form kinetic-energy matrix for a sinc (unbounded-domain) DVR basis \cite{colbertmiller1992},
\begin{equation}
T_{ij} = \frac{\hbar^2}{2m\Delta x^2}\times
\begin{cases}
\dfrac{\pi^2}{3}, & i=j,\\[4pt]
\dfrac{2\,(-1)^{i-j}}{(i-j)^2}, & i\neq j.
\end{cases}
\label{eq:kinetic}
\end{equation}
This formula is exact, not an approximation, for any wavefunction whose momentum content stays within the grid's Nyquist limit (Section~\ref{sec:dvr}-A) -- which is exactly the condition the grid spacing is chosen to guarantee. The diagonal entries cannot simply be read off the $i\neq j$ formula by setting $i=j$: doing so divides by zero rather than approaching any sensible limit, so the diagonal case is worked out on its own and turns out to be a different, finite expression. Because the off-diagonal formula depends only on how far apart two grid points are, $i-j$, and not on their absolute positions, the matrix's diagonals each repeat the same number over and over -- so only $N$ distinct values ever need to be computed, one per possible spacing, instead of $N^2$, one per matrix entry; the rest of the matrix is simply those $N$ numbers copied into place (Toeplitz matrix). Diagonalizing this matrix together with the potential then only needs to return the lowest $N_{\text{lev}}$ eigenvalues, not all $N$ of them; the linear-algebra routine used here (\texttt{eigvalsh} with an explicit index range) can extract just that low-energy slice directly, without the wasted work of also computing the high-energy states nobody asked for (LAPACK's MRRR algorithm).

\subsection{Why Only Smooth Potentials}
The sampling argument of Section~\ref{sec:dvr}-A also explains why the solver is restricted to potentials with no true discontinuity. At a hard wall the wavefunction's derivative jumps, its momentum content becomes effectively unbounded, and no finite grid spacing can resolve it -- attempting to do so aliases the unresolvable high-momentum content back into the low-lying, physically relevant part of the spectrum. Equation~\eqref{eq:kinetic} itself is also only valid where $V(x)$ is finite everywhere the grid is evaluated; the implementation raises an explicit error rather than silently return a wrong answer if a potential violates this.

\subsection{Automatic Grid Configuration}
The grid span and spacing are set automatically from the potential itself, rather than by hand. An energy ceiling $E_{\text{ceil}}\approx v_{\min}+1.5\,N_{\text{lev}}$ -- calibrated for the roughly-linear level spacing of an HO-like well, and only ever used as a starting estimate -- fixes a target grid extent via the classical turning points where $V(x)=E_{\text{ceil}}$, located by root-finding; the grid is padded modestly beyond them. The spacing is then set from the local wavenumber at that ceiling, $k_{\max}=\sqrt{2m(E_{\text{ceil}}-v_{\min})}/\hbar$ (derived in Appendix~\ref{sec:derivations}-D), via $\Delta x = \pi/(2k_{\max})$ -- twice as fine as the bare Nyquist limit $\pi/k_{\max}$, a safety margin against the initial estimate being slightly optimistic. Because this first estimate is only a starting point, an adaptive loop then widens the span geometrically (holding $\Delta x$ fixed) and re-solves, comparing eigenvalues against the previous, narrower grid; once the change falls below a set tolerance, the span is declared converged. This corrects for anharmonic or multi-well potentials -- the double well among them -- whose wavefunction tails extend well beyond what the HO-calibrated initial estimate predicts.

\subsection{Limitations of the Method}
Three limitations follow directly from the construction above, and bound what the pipeline can be trusted to produce. First, only roughly the lower half of any requested spectrum is numerically trustworthy: the highest-index states are the ones whose local momentum first exceeds what a fixed grid spacing can resolve, so $N_{\text{lev}}$ must always be requested with headroom beyond the levels actually needed (Section~\ref{sec:systems} quantifies this directly). Second, the method is restricted to smooth, single-valued, everywhere-finite potentials by Section~\ref{sec:dvr}-C; genuinely discontinuous potentials require a different DVR formulation entirely. Third, the dense $N\times N$ diagonalization means both memory and runtime scale unfavourably with the number of levels requested ($\mathcal{O}(N^2)$ and $\mathcal{O}(N^3)$ respectively) -- the practical ceiling on how far $N_{\text{lev}}$ can be pushed for a harder, anharmonic spectrum, and the motivation for the FFT-based alternative of Section~\ref{sec:future}.

\begin{figure*}[!t]
\centering
\includegraphics[draft=false,width=0.85\textwidth]{fig_pipeline.png}
\caption{Numerical workflow: solve the base spectrum, generate a higher-resolution numerical reference, verify convergence against it, compute quantum $C_v(T)$ and locate the classical plateau via the $\xi/n$-scan, and benchmark base- against reference-grid $C_v(T)$. The grid is refined automatically if the convergence check fails.}
\label{fig:pipeline}
\end{figure*}

% ===================== IV. NUMERICAL TECHNIQUES =====================
\section{Numerical Techniques and Implementation Safeguards}
\label{sec:tricks}
Beyond the grid construction of Section~\ref{sec:dvr}, several implementation-level techniques recur throughout the codebase, each addressing a specific way a naive implementation would silently fail or become impractically slow.

\textbf{Boltzmann-weight exponent shift.} Every partition-function sum in the pipeline -- the quantum $C_v(T)$, and every step of the $\xi$- and $n$-scans -- is evaluated as $w_n=\exp(-a_n+\min_n a_n)$ rather than the naive $\exp(-a_n)$, where $a_n=\beta E_n$ (or its $\xi$-scaled form). Subtracting the minimum exponent before exponentiating leaves every ratio $w_n/Z$ exactly unchanged, since it is common to every term, while preventing floating-point overflow or underflow when $a_n$ spans a wide range -- which it routinely does deep in the $\xi$-scan, where $\beta E_{\max}/\xi^2$ can be large. This is the standard log-sum-exp stability trick, applied here to every Boltzmann weight the pipeline computes.

\textbf{Vectorized $n$-convergence.} Rather than recomputing the partition sum from scratch for every candidate level count $n=2,\dots,N_{\text{lev}}$, the $n$-convergence scan of Section~\ref{sec:validation} builds the weighted sums $\sum_{k\le n}w_k$, $\sum_{k\le n}w_kE_k$, and $\sum_{k\le n}w_kE_k^2$ as a single running (cumulative) sum, giving $C_v(n)$ for every $n$ in one $\mathcal{O}(N_{\text{lev}})$ pass instead of $\mathcal{O}(N_{\text{lev}}^2)$.

\textbf{Partial diagonalization.} As noted in Section~\ref{sec:kinetic}, only the lowest $N_{\text{lev}}$ eigenvalues of $H$ are ever extracted, via an eigensolver that accepts an explicit index range rather than diagonalizing the full $N\times N$ matrix and discarding the unneeded upper eigenvalues -- a substantial saving whenever the grid (set by resolution, Section~\ref{sec:dvr}-D) is much larger than the number of levels actually requested.

\textbf{Decoupled span and resolution.} During the adaptive grid search of Section~\ref{sec:dvr}-D, the spacing $\Delta x$ is held fixed while the span is widened, so that a span-convergence failure and a resolution-convergence failure can never be conflated with each other -- each of the two error sources identified in Section~\ref{sec:dvr}-D is varied one at a time, never both together, which is what makes the two-pass convergence check of Section~\ref{sec:validation} a clean diagnostic rather than an ambiguous one.

% ===================== V. VALIDATION =====================
\section{Validation Strategy: A Ground Truth Without Closed-Form Solutions}
\label{sec:validation}
The pipeline is a short chain of system-agnostic stages (Fig.~\ref{fig:pipeline}): grid configuration and DVR solve, a higher-resolution numerical reference solve, a convergence check between the two, the $\xi/n$-scan, and a final base-vs-reference $C_v(T)$ benchmark. Switching to a new potential only requires changing the potential function itself; every other stage regenerates automatically.

\subsection{A Self-Consistent Reference, Not an Analytic One}
The ground truth at every stage is not a formula but a second DVR solve on an independently wider and finer grid. If the two agree to a set tolerance, the base result is declared converged -- the logic of Richardson extrapolation, used here as a pass/fail check rather than to improve the answer. This works identically for the HO, the double well, and every future potential, since it never touches a closed-form solution. It is not, however, proof of correctness by itself: a systematic error shared by both grids -- a bug in the DVR solver -- could make the two agree while both are wrong. Internal agreement rules out a coarse or truncated grid; it cannot rule out a flaw common to every grid the same code produces.

\subsection{Certifying the Method Against the One Known Answer}
That gap is closed by running the identical base-vs-reference machinery on the HO, without feeding it the Einstein formula of Eq.~\eqref{eq:cvho} anywhere, and only afterward comparing the result against it, external to the pipeline. Section~\ref{sec:systems} reports agreement to machine precision at every temperature: the self-consistent procedure converges to the physically correct answer, not merely to a shared artifact of the method -- the evidence that licenses trusting the same reference-grid check once the pipeline is pointed at potentials with no closed-form answer to compare against.

\subsection{Recognizing a Genuine Result vs.\ an Artifact}
A second check concerns the $\xi$-scan itself, and is where the two senses of ``Schottky anomaly'' in this report must be told apart. As $\xi$ increases, $C_v(\xi)$ is expected to plateau inside the window of Eq.~\eqref{eq:xiwindow}; pushed further, it instead collapses toward zero once the scaled spectrum runs out of thermally accessible levels. That collapse can reproduce the same rise-and-fall shape as the genuine, physical Schottky anomaly of Section~\ref{sec:intro} -- a finite spectrum, whatever its origin, is capable of the same signature -- so mistaking one for the other would mean reporting a spurious enhancement instead of the real effect being searched for. The pipeline tells them apart with a companion scan over how many computed levels were actually needed to reach the plateau: convergence using only a handful of low-lying levels reflects real physics, while a result that only stabilizes near the top of the computed spectrum is a warning that truncation is shaping the curve. Only a result clearing both checks -- base-vs-reference agreement \emph{and} a genuine, non-collapsing plateau -- is reported as the classical limit or as a candidate physical Schottky feature.

% ===================== VI. SYSTEMS AND RESULTS =====================
\section{Systems Studied and Results}
\label{sec:systems}

\subsection{Harmonic Oscillator -- Validation Case}
The HO is the system used to certify the methodology of Section~\ref{sec:validation}, since it is a simple case with an independent, closed-form answer to check against. Base-grid energy levels agree with the numerical reference to within machine precision across the full spectrum (Fig.~\ref{fig:ho_energy_comparison}), settling to a flat, floating-point-limited error plateau (Fig.~\ref{fig:ho_error_vs_n}). The resulting quantum and classical-limit $C_v(T)$ curves, together with the converged $\xi$ and $n$ values certifying each point, show the same agreement (Fig.~\ref{fig:cv_summary_ho}). Figures~\ref{fig:xi_convergence}--\ref{fig:n_convergence} show a representative pair of scans at the hardest temperatures in the sweep: a genuine $\xi$-plateau followed by the finite-$N$ collapse of Section~\ref{sec:validation} (Fig.~\ref{fig:xi_convergence}), and the companion $n$-convergence check confirming that plateau is reached using only a fraction of the computed spectrum (Fig.~\ref{fig:n_convergence}) -- the two outcomes and the two checks that Section~\ref{sec:validation}-C relies on.

\begin{figure*}[!t]
\centering
\includegraphics[draft=false,width=0.85\textwidth]{fig_ho_energy_comparison.png}
\caption{Base-grid vs. numerical-reference energy levels $E_n$ for the HO, full spectrum (left) and zoomed to the states with the largest absolute error (right).}
\label{fig:ho_energy_comparison}
\end{figure*}

\begin{figure}[!t]
\centering
\includegraphics[draft=false,width=\columnwidth]{fig_ho_error_vs_n.png}
\caption{Relative error between base-grid and reference-grid energy levels vs. state index $n$, settling to a flat machine-precision plateau -- consistent with both grids being converged and only floating-point round-off remaining.}
\label{fig:ho_error_vs_n}
\end{figure}

\begin{figure*}[!t]
\centering
\includegraphics[draft=false,width=0.7\textwidth]{fig_cv_summary_ho.png}
\caption{Quantum $C_v(T)$ (solid) and numerical classical limit (dashed) for the HO, with converged $\xi$ and $n$ on the secondary axis. The classical limit tracks the exact equipartition value $C_v/k_B=1$.}
\label{fig:cv_summary_ho}
\end{figure*}

The dashed plateau in Fig.~\ref{fig:cv_summary_ho} is not merely a numerical asymptote: it is the equipartition-theorem prediction for the HO's two quadratic degrees of freedom (kinetic and potential), $C_v=k_B$, derived independently of the $\xi$-scan in Appendix~\ref{sec:derivations}-C. That the numerically located plateau lands on this value, at every temperature in the sweep rather than only asymptotically at high $T$, is expected: for the HO the classical $C_v$ is exactly temperature-independent (Appendix~\ref{sec:derivations}-C), so the ``high-temperature'' equipartition value and ``the classical limit at this temperature'' coincide everywhere -- a correspondence that need not hold as cleanly for an anharmonic potential such as the double well (Section~\ref{sec:systems}-B).

\begin{figure}[!t]
\centering
\includegraphics[draft=false,width=\columnwidth]{fig_xi_convergence.png}
\caption{$\xi$-convergence at the hardest (lowest) temperature in the sweep: $C_v(\xi)$ rises (grey), stabilizes onto a genuine plateau (green), then collapses (yellow) once $\xi$ leaves the window of Eq.~\eqref{eq:xiwindow} -- the failure mode of Section~\ref{sec:validation}.}
\label{fig:xi_convergence}
\end{figure}

\begin{figure}[!t]
\centering
\includegraphics[draft=false,width=\columnwidth]{fig_n_convergence.png}
\caption{$n$-convergence at the hardest temperature in the sweep: $C_v(n)$ rises smoothly and converges using well under the full computed spectrum, confirming the plateau is not an artifact of the truncation.}
\label{fig:n_convergence}
\end{figure}

Separately, the DVR solver's own resolution and level-count limits were characterized directly (Figs.~\ref{fig:dvr_res_dx}--\ref{fig:dvr_res_n}), confirming the standard rule of thumb that only roughly the lower half of a computed spectrum should be trusted. Figures~\ref{fig:cv_benchmark_quantum}--\ref{fig:cv_benchmark_classical} close the loop with the external certification of Section~\ref{sec:validation}-B: the fully numerical, closed-form-free pipeline agrees with the exact Einstein-oscillator answer to machine precision, for both the quantum curve and its classical limit.

\begin{figure}[!t]
\centering
\includegraphics[draft=false,width=\columnwidth]{fig_dvr_res_dx.png}
\caption{DVR resolution limit: relative eigenvalue error vs. grid spacing $\Delta x$. Below a critical spacing the error stays at the machine-precision floor; beyond it, the error rises by many orders of magnitude.}
\label{fig:dvr_res_dx}
\end{figure}

\begin{figure}[!t]
\centering
\includegraphics[draft=false,width=\columnwidth]{fig_dvr_res_n.png}
\caption{DVR level-count limit: relative eigenvalue error vs. requested level index $n$, at fixed grid spacing. Only roughly the lower half of the computed spectrum remains trustworthy -- the highest-index states are the first whose local momentum exceeds what the grid can resolve.}
\label{fig:dvr_res_n}
\end{figure}

\begin{figure*}[!t]
\centering
\includegraphics[draft=false,width=0.8\textwidth]{fig_cv_benchmark_quantum_abs.png}
\caption{Quantum $C_v(T)$: fully numerical pipeline vs. the exact analytic Einstein-oscillator formula, with the absolute error shown below. Agreement is at machine precision throughout.}
\label{fig:cv_benchmark_quantum}
\end{figure*}

\begin{figure*}[!t]
\centering
\includegraphics[draft=false,width=0.8\textwidth]{fig_cv_benchmark_classical.png}
\caption{Classical-limit $C_v(T)$: fully numerical pipeline vs. the exact equipartition value $k_B$, with the relative error shown below. Every temperature point converges, at machine precision.}
\label{fig:cv_benchmark_classical}
\end{figure*}

\subsection{Double Well -- Next System}
The next system this project will tackle is a quartic double well, $V(x) = \tfrac14 x^4 + b\,x^3 - \tfrac12 x^2$, the first anharmonic potential in the pipeline. It is a harder test than the HO: its levels scale as $E_n\propto n^{4/3}$ rather than linearly, narrowing the plateau window of Eq.~\eqref{eq:xiwindow} (Section~\ref{sec:history}). The double well has no closed-form solution of its own, but it carries a validation route of a different, equally rigorous kind: near either minimum, any sufficiently smooth potential is locally well-approximated by a harmonic oscillator set by the local curvature, so the low-lying levels of a sufficiently deep well should closely match the already-certified HO benchmark of Section~\ref{sec:systems}-A -- an independent check without a second global closed-form answer. The double well is also the pipeline's first candidate for a genuine \emph{physical} Schottky anomaly: its central barrier supports a near-degenerate, tunnelling-split pair of levels well separated from the states above it, the textbook precondition for a real low-temperature $C_v(T)$ peak rather than a truncation artifact. This run is in progress; Section~\ref{sec:future} outlines the remaining steps.

% ===================== VII. VERSION HISTORY =====================
\section{Numerical Lessons This Semester}
\label{sec:history}
Manual grid tuning was replaced by the automatic Nyquist-based configuration of Section~\ref{sec:dvr}. A fixed padding heuristic, adequate for the HO, underestimated anharmonic wavefunction tails and was replaced with iterative span widening. Because the double well's $E_n\propto n^{4/3}$ scaling narrows its plateau window relative to the HO's linear spectrum, the $\xi$-scan step size was made more granular so the scan reliably lands inside the window before the finite-$N$ collapse of Section~\ref{sec:validation}. The DVR core was restricted to smooth potentials only, since a hard-wall discontinuity cannot be represented on any finite grid.

% ===================== VIII. FUTURE WORK =====================
\section{Future Work}
\label{sec:future}
The immediate priority for the continuation of the project is as follows: (i) complete and quantitatively characterize the double-well run, including the symmetric-vs-asymmetric comparison discussed with the supervisor; (ii) carry out the near-minimum harmonic-limit check of Section~\ref{sec:systems}-B explicitly, and test whether the tunnelling doublet produces a genuine, $n_{\text{conv}}$-verified Schottky peak; and (iii), motivated by Section~\ref{sec:intro}, scan systematically across potentials for a genuine Schottky-driven enhancement -- and, if one is found, quantify its efficiency against the classical limit, test whether it generalizes, and propose a concrete experimental realization.

Two structural changes to the solver were also identified (Section~\ref{sec:dvr}-E), but are worth pursuing only if that potential scan becomes compute-bound, rather than as ends in themselves: replacing the dense $N\times N$ diagonalization (which scales as $\mathcal{O}(N^3)$) with an FFT-accelerated, matrix-free Krylov-subspace eigensolver, since the kinetic operator is diagonal in momentum space; and an adaptive grid mapping that satisfies the Nyquist criterion locally rather than globally, expected to cut the required grid points by a factor of $3$--$4$.

% ===================== IX. CONCLUSION =====================
\section{Conclusion}
This semester's work built and validated the numerical groundwork for a longer-term search for quantum-enhanced heat machine performance: a pipeline that computes $C_v(T)$, quantum and classical, for any smooth potential, without depending on a closed-form solution for its own internal ground truth. That methodology was certified, once, against the HO, and reproduces it to machine precision -- evidence it can be trusted once no such answer exists to compare against. The same pipeline is now running, unmodified beyond the potential definition, on the anharmonic double well, whose near-minimum harmonic limit gives a further consistency check and whose tunnelling-split doublet is the first candidate for a genuine physical Schottky anomaly. Distinguishing that genuine enhancement from a numerical look-alike is exactly the distinction this project's longer-term goal depends on, and it is what this semester's validation machinery is designed to do.

\appendix
\section{Derivations}
\label{sec:derivations}

\subsection{Heat Capacity from the Partition Function (Eq.~\ref{eq:cv_general})}
For a canonical ensemble at inverse temperature $\beta$ with partition function $Z=\sum_n e^{-\beta E_n}$, the mean energy is $\expval{E}=-\partial_\beta \ln Z$. Differentiating $\ln Z$ directly,
\begin{equation}
\expval{E} = -\frac{1}{Z}\frac{\partial Z}{\partial\beta} = \frac{1}{Z}\sum_n E_n e^{-\beta E_n},
\end{equation}
which is the first moment of the energy under the Boltzmann weights. The heat capacity is $C_v=\partial_T\expval{E}=\left(\partial_\beta\expval{E}\right)\left(\partial_T\beta\right)$, and since $\beta=1/(k_BT)$ gives $\partial_T\beta=-k_B\beta^2$, only $\partial_\beta\expval{E}$ remains to be evaluated:
\begin{equation}
\frac{\partial\expval{E}}{\partial\beta} = -\frac{1}{Z^2}\frac{\partial Z}{\partial\beta}\sum_n E_ne^{-\beta E_n} - \frac{1}{Z}\sum_n E_n^2 e^{-\beta E_n} = \expval{E}^2-\expval{E^2}.
\end{equation}
Combining the two results, $C_v = -k_B\beta^2\left(\expval{E}^2-\expval{E^2}\right) = k_B\beta^2\,\mathrm{Var}(E)$, exactly Eq.~\eqref{eq:cv_general}. Nothing in this derivation refers to how $\{E_n\}$ was obtained, which is why the same formula applies unchanged to a closed-form spectrum or a numerically diagonalized one.

\subsection{The Valid Range for $\xi$ (Eq.~\ref{eq:xiwindow})}
The pipeline evaluates the scaled partition function $Z(\beta,\xi)=\sum_n e^{-\beta E_n/\xi^2}$. Two constraints on $\xi$ must hold simultaneously for the resulting $C_v(\xi)$ to represent the classical limit rather than an artifact.

\emph{Lower bound.} For the discrete sum to approximate the continuous classical integral it is meant to mimic, the exponent must change by much less than unity between adjacent levels, so that no single term dominates the local behaviour of the sum: with $\Delta E=E_{n+1}-E_n$ the local spacing, $\beta\Delta E/\xi^2\ll1$, i.e.
\begin{equation}
\xi \gg \sqrt{\beta\,\Delta E}.
\end{equation}
Below this, the spectrum still looks discrete to the scaled sum, and $C_v(\xi)$ is just the quantum $C_v$ again.

\emph{Upper bound.} Because the sum is truncated at a finite highest level $E_{\max}$, the truncation is only harmless if that level's Boltzmann weight stays negligible -- not that it "dominates" the sum, but that it must remain small enough not to distort the result relative to the true, untruncated sum: $e^{-\beta E_{\max}/\xi^2}\approx0$ requires the exponent to be large, $\beta E_{\max}/\xi^2\gg1$, i.e.
\begin{equation}
\xi \ll \sqrt{\beta\,E_{\max}}.
\end{equation}
Above this, the top of the truncated spectrum gains non-negligible weight and $C_v(\xi)$ collapses toward zero -- the finite-$N$ artifact of Section~\ref{sec:validation}. Combining the two gives the window of Eq.~\eqref{eq:xiwindow}, which only admits a valid $\xi$ at all when $\Delta E\ll E_{\max}$, i.e. when many levels are available.

\subsection{The Harmonic-Oscillator Heat Capacity (Eq.~\ref{eq:cvho}) and Its Classical Limit}
With $E_n=\hbar\omega(n+\tfrac12)$, the partition function is a geometric series,
\begin{equation}
Z^{\text{HO}} = e^{-\beta\hbar\omega/2}\sum_{n=0}^{\infty}\left(e^{-\beta\hbar\omega}\right)^n = \frac{e^{-\beta\hbar\omega/2}}{1-e^{-\beta\hbar\omega}},
\end{equation}
convergent since $e^{-\beta\hbar\omega}<1$ for $\beta,\hbar,\omega>0$. Differentiating $\expval{E}=-\partial_\beta\ln Z^{\text{HO}}$ gives $\expval{E}=\tfrac{\hbar\omega}{2}+\tfrac{\hbar\omega}{e^{\beta\hbar\omega}-1}$, and one further $\beta$-derivative (Appendix~\ref{sec:derivations}-A) gives, after using $\cosh$/$\sinh$ identities to simplify,
\begin{equation}
C_v^{\text{HO}} = k_B(\beta\hbar\omega)^2\frac{e^{\beta\hbar\omega}}{\left(e^{\beta\hbar\omega}-1\right)^2} = \frac14 k_B\left[\beta\hbar\omega\,\operatorname{csch}\!\left(\frac{\beta\hbar\omega}{2}\right)\right]^2,
\end{equation}
the two forms being algebraically identical since $\left(e^{x}-1\right)^2=e^{x}\left(e^{x/2}-e^{-x/2}\right)^2$.

Two independent routes confirm the classical limit is $k_B$. Taking $\beta\hbar\omega\to0$ in the exact formula above (L'H\^opital's rule on the indeterminate $0/0$ form) gives $C_v^{\text{HO}}\to k_B$. Separately, and without reference to the quantum formula at all, the equipartition theorem assigns $\tfrac12 k_BT$ to each quadratic term in the Hamiltonian $H=\tfrac{p^2}{2m}+\tfrac12 m\omega^2x^2$; with two such terms (kinetic and potential), $\expval{E}=k_BT$ and $C_v=\partial_T\expval{E}=k_B$, independent of $T$. That both routes -- one a limit of the discrete quantum spectrum, the other a classical phase-space argument that never mentions $\hbar$ -- agree exactly is the correspondence principle at work, and is what Fig.~\ref{fig:cv_summary_ho} confirms numerically.

\subsection{The Local Wavenumber $k_{\max}$ (Sections~\ref{sec:dvr}-A, D)}
For a state of energy $E$ in potential $V(x)$, classical energy conservation gives the local kinetic energy $p(x)^2/(2m)=E-V(x)$, so the local classical momentum is $p(x)=\sqrt{2m(E-V(x))}$. The de Broglie relation $p=\hbar k$ turns this into a local wavenumber $k(x)=p(x)/\hbar$ and a local wavelength $\lambda(x)=2\pi/k(x)$ that shrinks wherever the particle moves fastest (deep in the well) and grows toward the classical turning points, where $p(x)\to0$. Section~\ref{sec:dvr}-A requires the grid to resolve the \emph{shortest} such wavelength anywhere the wavefunction has appreciable amplitude, i.e. $p(x)$ maximized over $x$ at fixed $E$ -- which, from the formula above, occurs where $V(x)$ is smallest. Evaluating this at the potential minimum $v_{\min}$ for the highest requested energy $E_{\text{ceil}}$ gives exactly the $k_{\max}$ used to set the grid spacing in Section~\ref{sec:dvr}-D,
\begin{equation}
k_{\max} = \frac{\sqrt{2m\left(E_{\text{ceil}}-v_{\min}\right)}}{\hbar}.
\end{equation}

\section*{Code Availability}
The full source code, module-level documentation, figures, and reference literature cited in this report are maintained at:\\
\url{https://github.com/NadavJ180/quantum-classical-heat-capacity}

% ===================== REFERENCES =====================
\begin{thebibliography}{9}

\bibitem{gelbwaser2018}
D. Gelbwaser-Klimovsky, A. Bylinskii, D. Gangloff, R. Islam, A. Aspuru-Guzik, and V. Vuletic, ``Single-atom heat machines enabled by energy quantization,'' \emph{Phys. Rev. Lett.}, vol. 120, p. 170601, Apr. 2018.

\bibitem{colbertmiller1992}
D. T. Colbert and W. H. Miller, ``A novel discrete variable representation for quantum mechanical reactive scattering via the S-matrix Kohn method,'' \emph{J. Chem. Phys.}, vol. 96, no. 3, pp. 1982--1991, 1992.

\bibitem{lightcarrington2000}
J. C. Light and T. Carrington Jr., ``Discrete-variable representations and their utilization,'' in \emph{Advances in Chemical Physics}, vol. 114. Hoboken, NJ: John Wiley \& Sons, 2000, pp. 263--310.

\bibitem{campisifazio2016}
M. Campisi and R. Fazio, ``The power of a critical heat engine,'' \emph{Nat. Commun.}, vol. 7, p. 11895, 2016.

\bibitem{ye2017}
Z. Ye, Y. Hu, J. He, and J. Wang, ``Universality of maximum-work efficiency of a cyclic heat engine based on a finite system of ultracold atoms,'' \emph{Sci. Rep.}, vol. 7, p. 6289, 2017.

\bibitem{kozin2025}
V. K. Kozin, D. Miserev, D. Loss, and J. Klinovaja, ``Schottky anomaly in a cavity-coupled double quantum well,'' \emph{Phys. Rev. Res.}, vol. 7, p. 033239, 2025.

\bibitem{nyquist1928}
H. Nyquist, ``Certain topics in telegraph transmission theory,'' \emph{Trans. AIEE}, vol. 47, no. 2, pp. 617--644, Apr. 1928.

\bibitem{shannon1949}
C. E. Shannon, ``Communication in the presence of noise,'' \emph{Proc. IRE}, vol. 37, no. 1, pp. 10--21, Jan. 1949.

\end{thebibliography}

\end{document}
